\section{Future Work and Composability}
\label{storm:future-work}

Because we prove the generic construction secure with respect to our game-based
notions,
we do not inherently prove that the generic construction is secure with respect
to any other protocol;
i.e., we do not prove security under composability.
One approach to demonstrating composability is using the Universal Composability (UC) framework~\cite{Canetti01}.
However, a second approach would be to prove that a DKG that is secure under
these game-based notions is likewise secure under a simulation-based notion.
Such an approach was taken by Boneh et al.~\cite{BonehHL24} to show equivalence between
their game-based notions of security for an Exponent Verifiable Random Function (eVRF),
and a simulation-based notion of security.

To do so, we first must determine which notion of simulation-based security our game-based notions capture.
Full security defined a DKG whose output keys are indistinguishable from keys
from the target single-party key generation algorithm,
whereas security with abort defines a DKG that achieves full security with the
exception of an adversary that can force honest participants to abort (and
hence introduce bias).
